% ============================================================
% f2014_liaoning_science_q13_flowchart.tex
% 2014 辽宁卷（理科）第 13 题：程序框图
% 来源 PDF：普通高考/2014/2014辽宁理.pdf
% 原页：第 1 页（A3），印面 895，第 13 题题干下方
% 按原图重绘：输入 x；计算 y=x/3+2；判断 |y-x|<1，
% 成立则输出 y，不成立则令 x=y 后返回重新计算 y。
% ============================================================
\begin{tikzpicture}[>=Stealth, line width=0.7pt, font=\small,
  every node/.style={inner sep=2pt},
  term/.style={draw, rounded corners=3pt, minimum width=2.1cm,
    minimum height=0.6cm, align=center},
  proc/.style={draw, minimum width=2.5cm, minimum height=0.7cm, align=center},
  io/.style={draw, trapezium, trapezium left angle=70, trapezium right angle=110,
    minimum width=2.5cm, minimum height=0.6cm, align=center},
  dec/.style={draw, diamond, aspect=1.8, minimum width=2.7cm,
    minimum height=1.0cm, align=center}]
  \node[term] (start) at (0,0) {开始};
  \node[io] (input) at (0,-0.95) {输入 $x$};
  \node[proc] (calc) at (0,-1.95) {$y=\dfrac{x}{3}+2$};
  \node[dec] (test) at (0,-3.25) {$|y-x|<1$};
  \node[io] (out) at (0,-4.55) {输出 $y$};
  \node[term] (end) at (0,-5.5) {结束};
  \node[proc] (assign) at (3.2,-1.95) {$x=y$};

  \draw[->] (start) -- (input);
  \draw[->] (input) -- (calc);
  \draw[->] (calc) -- (test);
  \draw[->] (test.south) -- node[left] {是} (out.north);
  \draw[->] (out) -- (end);
  \draw[->] (test.east) -| node[pos=0.25, above] {否} (assign.south);
  \draw[->] (assign.north) |- (0,-1.45);
\end{tikzpicture}
